\section{A Gaussian residual operator bound}
\label{sec:residual}
Fix any probability measure $\xi$ on the input cube. It need not be known to the learner. All Hilbert norms in this section are in $L^2(\xi)$.

For an actual hidden layer, let
\[
 Q(x,y)=N^{-1}\sum_{i=1}^Nh_i(x)h_i(y),
 \qquad T c=N^{-1/2}\sum_i c_i h_i.
\]
Then $Q=TT^*$ as an integral operator and $\operatorname{tr}Q\le1$. At a spectral cutoff $\eta$, retain eigenvalues at least $\eta$, with factors $F_j=\sqrt{\lambda_j}e_j$. There are at most $1/\eta$ of them. The corresponding right-singular vectors $u_j\in\R^N$ are orthonormal and satisfy $T u_j=F_j$. Thus
\[
 F_j=N^{-1}\sum_i b_{ij}h_i,\qquad b_{ij}=\sqrt N u_j(i),
 \qquad N^{-1}\sum_i b_ib_i^\top=I.
\]
The residual kernel $E=Q-\sum_jF_j\otimes F_j$ is positive, has operator norm at most $\eta$, and has pointwise variance
$v(x)=Q(x,x)-\|F(x)\|^2\in[0,1]$.

Conditional on this preceding layer, the next Gaussian preactivation rows have the decomposition
\begin{equation}
 Z_i(x)=a_i\cdot F(x)+Z_i^0(x),
 \quad a_i\sim N(0,\tau^2I),\quad Z_i^0\sim\mathrm{GP}(0,\tau^2E),
 \label{eq:direct-decomposition}
\end{equation}
independently over rows and between the two parts. Define
\[
 m_i(x)=\E[\phi(Z_i(x))\mid a_i],\qquad
 r_i(x)=\phi(Z_i(x))-m_i(x),\qquad
 R c=N^{-1/2}\sum_i c_i r_i.
\]

\begin{lemma}[Residual feature concentration]
\label{lem:direct-residual}
For $K\ge1$, conditional on the preceding layer and every value of the latent vectors $(a_i)$, with probability at least $1-2e^{-NK/\tau^2}$,
\begin{align}
 \|R\|_\op&\le C[\omega^{-1/2}+\sqrt{\eta(K+1)}],\notag\\
 \left\|N^{-1}\sum_i\left(\phi(Z_i)^2-\E[\phi(Z_i)^2\mid a_i]\right)\right\|_{L^2(\xi)}
 &\le C[\omega^{-1/2}+\sqrt{\eta K}].
 \label{eq:direct-residual}
\end{align}
\end{lemma}
\begin{proof}
Fix a unit vector $c\in\R^N$. Conditional residual rows are independent and centered, and Gaussian Poincar\'e gives conditional pointwise variance at most $\tau^2v(x)$ because $\phi$ is 1-Lipschitz. Thus each residual row has integrated second moment at most $\tau^2$, and $\E\|Rc\|\le\tau/\sqrt N=\omega^{-1/2}$. As a function of all residual Gaussian coordinates, $\|Rc\|$ is Lipschitz with constant at most $\tau\sqrt{\eta/N}$. To verify this, dualize the Hilbert norm. The derivative of $\langle f,Rc\rangle$ in row $i$ has squared Gaussian-coordinate norm
\[
 \frac{c_i^2\tau^2}{N}
 \langle f\phi'(Z_i),E(f\phi'(Z_i))\rangle
 \le\frac{c_i^2\tau^2\eta}{N}\|f\|_2^2.
\]
Sum over rows and take $\|f\|_2=1$. Gaussian concentration therefore gives
\[
 \Pr(\|Rc\|>\omega^{-1/2}+t\mid(a_i))
 \le e^{-Nt^2/(2\tau^2\eta)}.
\]
A $1/2$-net of the unit sphere of $\R^N$ has at most $5^N$ points, and the operator norm is at most twice the maximum over this net. Choose $t=\sqrt{2\eta(K+\tau^2\log5)}$ and use a union bound. This proves the first assertion, uniformly over all coefficient vectors, including coefficients subsequently selected from the same residual realization.

For the second, Gaussian Poincar\'e gives the centered squared-feature Hilbert average expected norm at most $2\omega^{-1/2}$. Since $\phi^2$ is 2-Lipschitz, its norm is Gaussian-Lipschitz with constant at most $2\tau\sqrt{\eta/N}$. Concentration gives the stated bound directly. Infinite-dimensional Gaussian fields cause no additional issue: here every residual field is a linear image of finitely many row-weight coordinates, and the Hilbert-valued maps are Lipschitz.
\end{proof}

The standard Gaussian norm event
\begin{equation}
 \|W^1\|_F^2\le NK,\qquad
 \max_{2\le\ell\le H}\|W^\ell\|_\op^2\le K,\qquad
 N\|v\|^2\le K
 \label{eq:norm-event}
\end{equation}
has prior failure at most $(H+1)e^{-NK/(16\tau^2)}$ for a sufficiently large numerical lower bound on $K$. Its proof uses a chi-square tail for the first layer and readout and the Gaussian operator-norm tail for each internal matrix.

There is also an exact latent Gram bound on this event. If $V$ has the retained preceding right-singular vectors as columns, then $a_i=\sqrt N(WV)_{i,*}$, so
\begin{equation}
 N^{-1}\sum_i a_ia_i^\top=V^\top W^\top WV\preceq K I.
 \label{eq:actual-latent-gram}
\end{equation}
No independence between the retained current factors and their residuals is assumed. The sphere net in Lemma~\ref{lem:direct-residual} was taken precisely to control these adaptively selected coefficient vectors.
